\section{What is being counted?}

Fix an even natural number $N$. Both arguments seek a prime $p$ for which the
complementary integer $N-p$ has very few prime factors. Write $\Omega(n)$ for
the number of prime factors of a positive integer $n$, counted with
multiplicity. Thus a prime square has two prime factors for this purpose.
The Chen representation count in this project is
\[\begin{gathered}
 R_2(N)=\#\{p<N:p\text{ prime},\\
 N-p\ge2,\quad\Omega(N-p)\le2\}.
\end{gathered}\]
A member of this set gives a prime plus a prime or a product of two primes.
The two factors may coincide. The implementation names this finite set
\texttt{chenGoodRepresentations}; $R_2(N)$ is its cardinality.

Li--Liu impose a size condition on the two complementary factors. Their
count, written $D_{1,19/10}(N)$ here and \texttt{D19} in Lean, counts primes
$p\le N$ for which there exist natural numbers $r,q$ satisfying
\[\begin{gathered}
 N=p+rq,\qquad q\text{ prime},\\
 r=1\text{ or }r\text{ prime},\qquad r^{10}\le q^9.
\end{gathered}\]
Each eligible $p$ is counted once, regardless of the number of witnesses.
The exponent condition is stated in exact natural-number arithmetic; over
the nonnegative reals it is equivalent to $r\le q^{9/10}$.

For $N>1$, the common normalization is
\[
 \mathcal X_N=\mathfrak S_{\mathrm{Liu}}(N)\frac{N}{(\log N)^2},
\]
where
\[
 \mathfrak S_{\mathrm{Liu}}(N)=
 \prod_{\substack{\ell>2\ \mathrm{prime}\\\ell\mid N}}
       \frac{\ell-1}{\ell-2}
 \prod_{\ell>2\ \mathrm{prime}}\left(1-\frac{1}{(\ell-1)^2}\right).
\]
The proofs establish positivity of this scale for the relevant large $N$.
The public quantitative conclusions are $R_2(N)\ge0.67\mathcal X_N$ and
$D_{1,19/10}(N)>0.0004\mathcal X_N$, eventually for every even $N$.
The counting definitions and the proof routes remain separate.

\section{The mathematical plan}

Both proofs replace a difficult existence question by a finite counting
inequality. The finite inequality comes first: it determines precisely
which estimates the analytic part must supply. Positive contributions need
lower bounds, negative contributions need upper bounds, and all remainders
must fit into a single positive final margin.

For Chen, follow this route:
\begin{enumerate}
\item Build a weighted count of prime complements, together with penalties
      for the unwanted factor patterns. A finite counting bridge compares
      the corrected weight with the actual count $R_2(N)$.
\item Apply the weighted lower sieve to the original Goldbach source.
      The lower-bound construction includes its prime-power correction.
\item Bound the remaining triple-factor penalty through switching and the
      Selberg upper sieve. A distribution estimate for the switched source
      supplies the required averaged remainder control.
\item Put the lower and upper estimates into the finite bridge, absorb the
      errors in one eventual range, and obtain the representation bound.
      Positivity then yields a representation.
\end{enumerate}
Chapter~\ref{chap:chen} explains each of these steps and its exact implemented
interfaces. In particular, the triple-factor correction is a concrete
count, not a synonym for an unspecified error term.

For Li--Liu, follow a different finite reduction:
\begin{enumerate}
\item Construct weights whose pointwise inequalities detect the literal
      constrained-factor representations counted by $D_{1,19/10}(N)$.
\item Decompose the resulting weighted count into sifted sums and
      prime-factor ranges. Retain the signs and coefficients of the
      decomposition throughout the estimates.
\item Bound those terms with the appropriate lower or upper sieve, switched
      counts and prime-sum integrals. Certify the required bounds on the
      actual integrals and preserve the factor-range boundaries.
\item Combine the estimates with a common error budget. The implementation
      exposes both a positive-margin existence proof and a sharper
      quantitative proof using the author G11 estimate.
\end{enumerate}
Chapter~\ref{chap:liliu} defines the paper's numbered sums and integrals at
first use, before explaining how they enter these steps.

Chapter~\ref{chap:foundations} identifies the shared analytic inputs.
Those inputs are reused with the source, local density and remainder
weights appropriate to each application.

\section{Reading a sieve estimate}

The notation in this paragraph is a dictionary for the sieve interfaces.
Let $A$ be a finite set of integers and $\mathcal P$ a specified finite set
of sieving primes. Its sifted count is
\[
 S(A,\mathcal P)=\#\{a\in A:\ell\nmid a\text{ for every }\ell\in\mathcal P\}.
\]
In applications, weights or multiplicities replace the unweighted count.
For a positive divisor $d$, write $A_d=\{a\in A:d\mid a\}$ and separate
its count into an expected local-density term and a remainder:
\[
 |A_d|=g(d)X_A+r_d.
\]
Here $X_A$ is the source mass, $g(d)$ is its local density, and $r_d$ is the
error. The density is part of the source model; it is not interchangeable
between different sources. A sieve bound combines a main term built from
these densities with a weighted sum of the remainders.

A \emph{level of distribution} describes the divisor range over which a
suitable average of these errors can be controlled. A \emph{conditioned
source} additionally fixes one or more prime divisors before sieving the
remaining factor. \emph{Switching} reorganizes which factors are summed
first, so that the upper sieve can act on a more useful source. Each
application below specifies the actual carriers and cutoffs; this dictionary
does not replace their endpoint conventions.

\section{How to use the three layers}

Read the two argument outlines above for the global structure. A first-time
reader can go directly to either proof chapter and return to the analytic
foundations when an input is needed; the foundations need not be read straight
through before seeing an application. Cross-references next to the first uses
of numbered sieve terms lead to their exact definitions. The overview
graph compresses existing dependency paths between the main mathematical
stages; its arrows can stand for several intermediate lemmas. Then follow
the relevant proof chapter: its displayed inequalities explain what each
input proves and why the next step follows. Open a chapter's proof graph to
inspect the compiled dependencies between its named lemmas; nodes marked
\emph{input} are introduced in another chapter. Node titles are mathematical
names, and clicking a node opens its statement and exact Lean source.
On narrow screens, long displayed equations can be scrolled horizontally.

The all-declarations graph is also available for cross-chapter navigation.
Both graph views project onto explicitly documented lemmas. Their edges are
extracted from compiled declarations; the chapter prose explains the
mathematical reductions that a graph alone cannot express. Elementary Lean
plumbing and compatibility wrappers are left to the source view.
